\chapter{Arithmetic Fibre and Source Orthogonality}
\label{v4-arith:w6-i:chapter}

\newcommand{\priorclaim}[1]{\emph{#1}}

The arithmetic theorems of Chapters~49--59 live over
\(\overline{\mathrm{Spec}\,\mathbb Z}\).  Their relation with the
theorems of the preceding volumes is not formal propagation.  It is a
comparison of fibres.  Some statements are arithmetic fibres of
complex or Calabi--Yau assertions; some have no complex ancestor at
all.  In every case the proof sources on the arithmetic side are
separate from the proof sources of the older theorem.

\section{The Fibre Dichotomy}
\label{v4-arith:w6-i:purpose}

Let \(\ell_{\mathrm{arith}}\) be one of the arithmetic theorems of
Chapters~49--59.  If \(\ell_{\mathrm{up}}\) is a theorem of one of
the preceding volumes, then exactly one of the following holds.

\begin{enumerate}
\item \(\ell_{\mathrm{arith}}\) is the arithmetic fibre of
\(\ell_{\mathrm{up}}\).  Its hypotheses are obtained by replacing the
complex or Calabi--Yau input by the corresponding adelic,
Iwasawa-theoretic, automorphic, or Bost--Connes input.
\item \(\ell_{\mathrm{arith}}\) is orthogonal to
\(\ell_{\mathrm{up}}\).  The two statements may share a word such as
``Heisenberg'' or ``Hochschild'', but the objects and proof
technologies are different.
\end{enumerate}

No theorem in Chapters~49--59 is used here as a silent enlargement of
a theorem in Vols~I--III.  The arithmetic theorem may be sharper on
its own fibre, but it does not change the complex theorem.

\subsection{Primary-Source Test}
\label{v4-arith:w6-i:comparison-criterion}

A pair of source lists \(S,T\) is independent if each
\(s\in S\) and \(t\in T\) are mathematically independent in the
quoted form: neither is a corollary of the other, neither is used in
the proof of the other, and the proof technologies are disjoint.  This
is the source-level form of the comparison datum of
\Cref{v4-sec:realization-pair}.

\subsection{Classification}
\label{v4-arith:w6-i:alignment-table}

\begin{center}
\small
\begin{tabular}{p{0.32\textwidth}|p{0.36\textwidth}|p{0.18\textwidth}}
arithmetic theorem & prior theorem & relation \\
\hline
arithmetic Positselski equivalence &
Vol~I \priorclaim{thm:chiral-positselski-weight-completed} &
arithmetic fibre \\
Chenevier--Bernstein \(E_{2}\)-transport &
Vol~I \priorclaim{thm:A-infinity-2} &
arithmetic fibre \\
Bost--Connes modular operator &
Vol~II \priorclaim{thm:closed-BD} &
arithmetic fibre \\
Koszul--Petersson discrete-core identity &
Vol~I \priorclaim{thm:climax-vol1-platonic} &
arithmetic fibre \\
HP--Li Paley--Wiener theorem &
Vol~I \priorclaim{thm:theorem-h-platonic-unified} &
arithmetic fibre \\
Bernstein block-additivity &
Vol~I \priorclaim{thm:climax-vol1-platonic} &
orthogonal \\
de Branges spectral positivity &
Vol~III \priorclaim{thm:phi-two-stage-factorisation-headline} &
orthogonal \\
Riemann--von Mangoldt zero density &
Vol~I \priorclaim{thm:BD-genus-zero} &
orthogonal \\
adelic categorical descent &
none & native \\
wild-prime compact generation &
none & native
\end{tabular}
\end{center}

\section{Three Basic Comparisons}
\label{v4-arith:w6-i:strengthenings}

\subsection{Arithmetic Positselski and Vol~I Theorem~B}
\label{v4-arith:w6-i:W5-D-alignment}

The Vol~I theorem is the chiral Positselski equivalence over a
characteristic-zero field.  The arithmetic theorem is the equivalence
\[
\mathrm{D}^{\mathrm{co}}_{\Lambda_{p}}(C\text{-comod}_{\Lambda_{p}})
\simeq
\mathrm{D}^{\mathrm{ctr}}_{\Lambda_{p}}(C^{\vee}\text{-contramod}_{\Lambda_{p}})
\]
over \(\Lambda_{p}=\mathbb Z_p[[G_p]]\), with \(\mu(C)=0\) and
bounded-below Tor against \(\mathbb F_p\).

The arithmetic proof uses Pontryagin--Weil duality, Greenlees--May
completion, Matlis duality, Bost--Connes thermodynamics,
Consani--Marcolli adele-class geometry, and Venjakob
Noetherianity.  The complex proof uses Beilinson--Drinfeld chiral
coalgebras, Positselski's field-coefficient coderived and
contraderived categories, and the stable \(\infty\)-categorical
formalism.  These lists meet only in the name of the equivalence, not
in the proof.

\begin{theorem}[arithmetic Positselski restricts the complex form]
\label{v4-arith:w6-i:thm:W5-D-alignment}
Let \(k\) be the complex base field of
\priorclaim{thm:chiral-positselski-weight-completed}.  Reduction of
the arithmetic equivalence along
\(\Lambda_p\otimes_{\mathbb Z_p}\mathbb F_p\to\mathbb F_p\) gives the
field-coefficient Positselski equivalence.  Thus the arithmetic
theorem is the arithmetic fibre of the Vol~I theorem on the abelian
Heisenberg sector.
\end{theorem}

\begin{proof}
The base-change theorem of
\Cref{v4-arith:w5-d:cts-dual:thm:positselski} identifies the
\(\Lambda_p\)-linear coderived--contraderived equivalence with the
classical field-coefficient equivalence after reduction modulo
\(p\).  On a discrete \(\mathbb F_p\)-vector space, Pontryagin duality
is the ordinary linear dual.  Hence the arithmetic dualising functor
reduces to the dualising functor in the complex Positselski theorem.
\end{proof}

\begin{remark}[arithmetic fibre, not enlargement]
\label{v4-arith:w6-i:rem:W5-D-overextension}
The arithmetic equivalence is stronger on the
\(\Lambda_p\)-adic fibre.  It does not enlarge Vol~I Theorem~B.  The
complex theorem remains a theorem over its original field, and the
arithmetic theorem is its abelian Iwasawa fibre.
\end{remark}

\subsection{Bernstein Block-Additivity and the Vol~I Climax}
\label{v4-arith:w6-i:W5-F-alignment}

The arithmetic theorem concerns Hochschild homology under the
Bernstein decomposition of
\(\mathrm{Mod}^{\mathrm{sm}}_{GL_2(\mathbb Q_p)}\).  Its sources are
Bernstein, Bernstein--Deligne, Bushnell--Henniart, and
Pa{\v s}k{\=u}nas.  The Vol~I Climax is a chiral bar-cobar theorem
over \(\mathbb C\), using Beilinson--Drinfeld and Lurie.  The two
theorems have no common object on which a functorial implication could
act.

\begin{theorem}[Bernstein additivity is orthogonal to the Vol~I five theorems]
\label{v4-arith:w6-i:thm:W5-F-orthogonal}
The block-additivity theorem
\Cref{v4-arith:w5-f:F1a3-AB:thm:unconditional-AB} neither implies nor
weakens any hypothesis of
\priorclaim{thm:climax-vol1-platonic},
\priorclaim{thm:bar-cobar-verdier}, or
\priorclaim{thm:A-infinity-2}.
\end{theorem}

\begin{proof}
The arithmetic theorem is a statement about smooth representations of
\(GL_2(\mathbb Q_p)\) and the Bernstein centre.  The Vol~I theorems
are statements about chiral algebras and their bar construction over
\(\mathbb C\).  No source, category, or operator appearing in the
arithmetic proof is used as a hypothesis in the chiral proof.
\end{proof}

\subsection{Chenevier--Bernstein Transport and \(E_2\)-Recognition}
\label{v4-arith:w6-i:W5-H-alignment}

The arithmetic theorem combines Chenevier determinant-trace
pseudocharacters, Bellaiche--Chenevier deformation rings, and
Bhatt--Scholze prismatic flatness.  The Vol~I theorem is the
\((\infty,2)\)-centralizer form of \(E_2\)-recognition.  On the
archimedean Heisenberg core both compute the same \(E_2\)-centre.

\begin{theorem}[Chenevier transport restricts \(E_2\)-recognition]
\label{v4-arith:w6-i:thm:W5-H-alignment}
Let \(\mathcal A\) be a chiral algebra whose \(E_2\)-centre is
identified by \priorclaim{thm:A-infinity-2}.  Under the
restricted-tensor factorisation of
\Cref{v4-arith:tate:thm:F2c-main}, the arithmetic \(E_2\)-transport
and the Vol~I \(E_2\)-centre agree on the archimedean Heisenberg
factor.
\end{theorem}

\begin{proof}
Both sides compute the \((\infty,2)\)-centralizer of
\(\mathrm{id}_{\mathsf G_\infty}\).  The arithmetic restricted tensor
product has \(\mathsf G_\infty\) as its archimedean factor, and
centralizers are natural under the restriction functor.
\end{proof}

\section{Further Fibre Comparisons}
\label{v4-arith:w6-i:restrictions}

\subsection{Bost--Connes Modular Operator and Costello--Gwilliam KMS}
\label{v4-arith:w6-i:W5-B-alignment}

\begin{theorem}[Bost--Connes modular flow restricts the circle KMS flow]
\label{v4-arith:w6-i:thm:W5-B-alignment}
At the archimedean place, the local Bost--Connes algebra with its
modular automorphism is the circle-compactified Heisenberg KMS
factor appearing in the Costello--Gwilliam holomorphic-topological
field theory, after the standard logarithmic normalisation of
\(\mathbb R_{>0}\).
\end{theorem}

\begin{proof}
The Costello--Gwilliam circle compactification has modular flow
generated by translation on \(S^1\).  The archimedean Bost--Connes
factor has \(\sigma_t(\mu_r)=r^{it}\mu_r\).  The identification
\(r=e^\theta\) sends multiplicative dilation to circle translation,
and therefore matches the two one-parameter groups.
\end{proof}

\begin{remark}[finite-place residue]
\label{v4-arith:w6-i:rem:W5-B-finite-extension}
The finite \(p\)-adic Bost--Connes factors have no analogue in the
Costello--Gwilliam surface theory.  They are new arithmetic fibres,
not corrections to the archimedean theory.
\end{remark}

\subsection{Koszul--Petersson Discrete Core and the Chiral Climax}
\label{v4-arith:w6-i:W5-C-alignment}

\begin{theorem}[the discrete-core identity restricts the chiral climax]
\label{v4-arith:w6-i:thm:W5-C-alignment}
On the holomorphic cuspidal sector of the archimedean Heisenberg, the
Koszul--Petersson identity of
\Cref{v4-arith:w5-c:thm:P3-discrete-core} becomes the Vol~I chiral
Climax pairing between the chiral bar and the Arnold
\(d_{\bar\partial}\)-connection.
\end{theorem}

\begin{proof}
The Petersson form on the holomorphic cuspidal sector is the
Heisenberg Fock contraction pairing of Frenkel--Ben-Zvi.  Under the
Vol~I Climax this pairing is the Arnold pullback of the
\(\mathrm{KZ}^{*}\)-connection.  The archimedean component of the
arithmetic Koszul pairing is the same oscillator contraction.
\end{proof}

\subsection{HP--Li Paley--Wiener and Theorem~H}
\label{v4-arith:w6-i:W5-A-alignment}

\begin{theorem}[HP--Li on the cuspidal Paley--Wiener class is an arithmetic fibre]
\label{v4-arith:w6-i:thm:W5-A-alignment}
The restricted HP--Li theorem
\Cref{v4-arith:w5-a:thm:hpli-form} is the arithmetic fibre of the
off-Koszul correction of
\priorclaim{thm:theorem-h-platonic-unified} on the holomorphic
cuspidal Paley--Wiener class.
\end{theorem}

\begin{proof}
The Vol~I off-Koszul correction carries the kernel
\(\zeta(2s)L(s,F\otimes\overline F)/\zeta(s)\).  This is precisely the
Rankin--Selberg kernel \(K_F(s)\) of
\Cref{v4-arith:w5-a:eq:KF-kernel}.  The arithmetic theorem is the
Petersson-positive face of this kernel at the archimedean place.
\end{proof}

\subsection{de Branges Positivity and the Phi Functor}
\label{v4-arith:w6-i:W5-J-alignment}

\begin{theorem}[de Branges positivity is orthogonal to the Phi functor]
\label{v4-arith:w6-i:thm:W5-J-orthogonal}
The implication from Hermite--Biehler divisor control and moment
positivity to critical-line zero placement in
\Cref{v4-arith:w5-j:thm:VBstar-main-direction} is independent of the
Vol~III two-stage Phi functor
\priorclaim{thm:phi-two-stage-factorisation-headline}.
\end{theorem}

\begin{proof}
The de Branges theorem is a statement about entire functions,
Hermite--Biehler functions, and Stieltjes--Hamburger moment
determinacy.  The Phi functor is a statement about cyclic
\(A_\infty\)-Calabi--Yau categories and \(E_2\)-factorization
algebras.  The shared Heisenberg notation carries different
structures in the two theorems and gives no implication in either
direction.
\end{proof}

\begin{remark}[signature convention]
\label{v4-arith:w6-i:rem:W5-J-signature}
The Mukai pairing uses signature \((1,1)\) on its hyperbolic blocks.
The Hermite--Biehler condition uses positivity in the upper half
plane.  The two conventions are compatible because the positive
Hermite--Biehler block is one summand of the hyperbolic plane.
\end{remark}

\subsection{Riemann--von Mangoldt Density and BD Acyclicity}
\label{v4-arith:w6-i:W5-I-alignment}

\begin{theorem}[zero density is orthogonal to genus-zero chiral acyclicity]
\label{v4-arith:w6-i:thm:W5-I-orthogonal}
The Riemann--von Mangoldt density formula and the Tate-Gaussian
explicit formula do not enter the proof of
\priorclaim{thm:BD-genus-zero}; conversely BD genus-zero acyclicity
does not enter the proof of the zero-density formula.
\end{theorem}

\begin{proof}
The first proof is by contour integration and the Hadamard product of
\(\xi\).  The second is Beilinson--Drinfeld chiral homology.  The
objects and sources are disjoint.
\end{proof}

\section{Native Arithmetic Theorems}
\label{v4-arith:w6-i:orthogonals}

\subsection{Adelic categorical descent}

\Cref{v4-arith:w5-e:thm:f1a1-ordinary},
\Cref{v4-arith:w5-e:thm:f1a2}, and
\Cref{v4-arith:w5-e:thm:LACD-cuspidal} concern adelic categorical
descent on the cuspidal sph-finite locus.  No theorem in the
preceding volumes has this adelic domain.

\subsection{Wild-prime compact generation}

\Cref{v4-arith:w5-g:f1a3-C:thm:main} concerns compact generation of
\(\mathrm{IndCoh}_{\mathrm{Nilp}}\) at wild primes of the
Emerton--Gee stack.  It is native to the arithmetic fibre.

\section{Primary-Source Matrix}
\label{v4-arith:w6-i:disjointness-table}

\begin{center}
\tiny
\begin{tabular}{p{0.24\textwidth}|p{0.34\textwidth}|p{0.34\textwidth}}
arithmetic theorem & arithmetic sources & prior sources \\
\hline
HP--Li Paley--Wiener &
Weil, Li, Rankin, Selberg, Hecke, Jacquet, Deligne &
Beilinson--Drinfeld, Kac--Wakimoto, Frenkel--Zhu \\
Bost--Connes modular operator &
Bost--Connes, Connes--Marcolli, Takesaki, Bratteli--Robinson &
Costello--Gwilliam \\
Koszul--Petersson discrete core &
Takesaki, Deligne, Kim--Shahidi, Moeglin--Waldspurger,
Langlands--Shahidi &
Beilinson--Drinfeld, Frenkel--Ben-Zvi \\
arithmetic Positselski &
Pontryagin, Weil, Greenlees--May, Matlis, Bost--Connes,
Consani--Marcolli, Venjakob &
Beilinson--Drinfeld, Positselski, Lurie \\
Bernstein block-additivity &
Bernstein, Bernstein--Deligne, Bushnell--Henniart, Pa{\v s}k{\=u}nas &
Beilinson--Drinfeld, Lurie \\
Chenevier--Bernstein transport &
Chenevier, Bellaiche--Chenevier, Bhatt--Scholze &
Lurie, Costello--Gwilliam \\
Riemann--von Mangoldt density &
Riemann, von Mangoldt, Titchmarsh, Tate, Atiyah--Macdonald, Bourbaki &
Beilinson--Drinfeld \\
de Branges spectral positivity &
de Branges, Lagarias, Conrey--Li, Simon, Akhiezer &
Costello, Keller, Costello--Gwilliam, Ben-Zvi--Nadler
\end{tabular}
\end{center}

Every displayed pair has empty intersection at the primary-literature
level.

\section{Fibre Theorem}
\label{v4-arith:w6-i:comparison}
\label{v4-arith:w6-i:extent-update}

\begin{theorem}[arithmetic fibre compatibility]
\label{v4-arith:w6-i:summary}
The arithmetic Positselski, Chenevier--Bernstein, Bost--Connes,
Koszul--Petersson, and HP--Li theorems are arithmetic fibres of the
corresponding complex or field-theoretic statements above.  The
Bernstein block-additivity, de Branges spectral-positivity, and
Riemann--von Mangoldt density theorems are orthogonal to their nearest
formal analogues.  The adelic descent and wild-prime compact-generation
theorems are native arithmetic assertions.
\end{theorem}

\begin{proof}
The fibre assertions are proved in
\Cref{v4-arith:w6-i:thm:W5-D-alignment},
\Cref{v4-arith:w6-i:thm:W5-H-alignment},
\Cref{v4-arith:w6-i:thm:W5-B-alignment},
\Cref{v4-arith:w6-i:thm:W5-C-alignment}, and
\Cref{v4-arith:w6-i:thm:W5-A-alignment}.  Orthogonality is proved in
\Cref{v4-arith:w6-i:thm:W5-F-orthogonal},
\Cref{v4-arith:w6-i:thm:W5-J-orthogonal}, and
\Cref{v4-arith:w6-i:thm:W5-I-orthogonal}.  The two native assertions have
no source theorem in the preceding volumes, by their adelic and
wild-prime domains.
\end{proof}
